\answer
\begin{enumerate}

%% 6章 問1（6.1節）
\item[6.1節 問1]
\rev{$V_{\mathrm{out}} = DV_{\mathrm{in}}$より}
\begin{equation*}
\rev{D = \frac{V_{\mathrm{out}}}{V_{\mathrm{in}}} = \frac{5}{141} \approx 0.035}
\end{equation*}
\rev{周期は$T = 1/f = 10\,\mathrm{\mu s}$だから，オン時間は
$T_{\mathrm{on}} = DT \approx 0.35\,\mathrm{\mu s} = 350$\,nsしかない。}
スイッチの切替えに数十\,nsかかるとすれば，
オン時間のかなりの部分が切替えの途中の状態で占められ，
オン時間を正確に制御できないうえ，スイッチング損失の割合も大きくなる。
大比率の変換をデューティ比だけで受け持たせるのは現実的でない
（巻数比に受け持たせる絶縁型が有利になる）。

%% 6章 問2（6.2節）
\item[6.2節 問2]
式\ref{eq6.1}の右辺：$N_1$は無次元だから，単位は
$\dfrac{\mathrm{Wb}}{\mathrm{s}} = \dfrac{\mathrm{V\cdot s}}{\mathrm{s}}
= \mathrm{V}$となり，左辺（電圧）と一致する。
式\ref{eq6.4}の右辺：
$\mathrm{H}\times\dfrac{\mathrm{A}}{\mathrm{s}}
= \dfrac{\mathrm{V\cdot s}}{\mathrm{A}}\times\dfrac{\mathrm{A}}{\mathrm{s}}
= \mathrm{V}$となり，こちらも左辺と一致する。

%% 6章 問3（6.3節）
\item[6.3節 問3]
リセット期間には$\mathrm{D_3}$が導通し，リセット巻線が電源につながれるから，
1次巻線の電圧は式\ref{eq6.13}より
\rev{$v_1 = -({N_1}/{N_3})V_{\mathrm{in}} = -V_{\mathrm{in}}$（$N_3 = N_1$）である。}
スイッチの両端電圧は電源電圧と1次巻線電圧（の大きさ）の和だから
\begin{equation*}
\rev{v_{sw} = V_{\mathrm{in}} + \frac{N_1}{N_3}V_{\mathrm{in}} = 2V_{\mathrm{in}}}
\end{equation*}
となる。フォワードコンバータのスイッチには入力の2倍の電圧がかかるので，
\rev{たとえば$V_{\mathrm{in}} = 141$\,Vなら$282$\,Vに余裕を見た耐圧の素子を選ぶ必要がある。}

%% 6章 問4（6.4節）
\item[6.4節 問4]
\rev{$W = \frac{1}{2}L_m I_{\mathrm{pk}}^2$の右辺：}
$\mathrm{H}\times\mathrm{A^2}
= \dfrac{\mathrm{V\cdot s}}{\mathrm{A}}\times\mathrm{A^2}
= \mathrm{V\cdot A\cdot s} = \mathrm{W\cdot s} = \mathrm{J}$となり，
エネルギーの単位と一致する。
\rev{$P = Wf$の右辺：}
$\mathrm{J}\times\dfrac{1}{\mathrm{s}} = \mathrm{J/s} = \mathrm{W}$となり，
電力の単位と一致する。
「1回に運ぶエネルギー$\times$毎秒の回数＝電力」という
フライバックの動作そのものを表す検算である。

\end{enumerate}
